%%%%%%%%%%%%%%%%%%%%%%%%%%%%%%%%%%%%%%%%%%%%%%%%%%%%%%%%%%%%%%%%%%%%%%%%%%%%%%%
%% Aalto University MSc Thesis — AI-EDITED VERSION
%% Author: Aleksander Heino
%% Topic: Encoding-Surrogate Pairing Effects in ML-Guided Protein Optimisation
%% Build: latexmk -pdf thesis_ai_edited
%%%%%%%%%%%%%%%%%%%%%%%%%%%%%%%%%%%%%%%%%%%%%%%%%%%%%%%%%%%%%%%%%%%%%%%%%%%%%%%

\documentclass[english, 12pt, a4paper, elec, utf8, a-2b, online]{aaltothesis}
%\documentclass[english, 12pt, a4paper, elec, utf8, a-2b, print]{aaltothesis}

\setupthesisfonts
\usepackage{graphicx}
\usepackage{longtable}
\usepackage{booktabs}
\usepackage{amsmath}
\usepackage[type={CC}, modifier={by-nc-sa}, version={4.0}]{doclicense}

\usepackage[
  backend=biber,
  style=numeric-comp,
  sorting=nyt,
  firstinits=true,
  urldate=long
]{biblatex}
\addbibresource{thesisreferences.bib}

%% --- Fill in your details ---
\degreeprogram{Electronics and Electrical Engineering}
\major{Machine Learning and Data Science}  %% TODO: confirm your exact major name
\univdegree{MSc}
\thesisauthor{Aleksander Heino}
\thesistitle{Encoding--Surrogate Pairing Effects in Machine Learning-Guided Protein Fitness Optimisation}
\thesissubtitle{A Systematic Evaluation on Sparse and Dense Fitness Landscapes}
\place{Espoo}
\date{\today}
\supervisor{Prof.\ Supervisor Name}   %% TODO: fill in
\advisor{Dr Advisor Name}             %% TODO: fill in
\collaborativepartner{}
\uselogo{'}

\copyrighttext{\noexpand\textcopyright\ \number\year\ \ThesisAuthor. This work is
  licensed under a Creative Commons ``Attribution-NonCommercial-ShareAlike 4.0
  International'' (BY-NC-SA 4.0) license.}{\noindent\textcopyright\ \number\year\ \ThesisAuthor\ \doclicenseThis}

\keywords{protein engineering\spc machine learning\spc Bayesian optimisation\spc protein language models\spc fitness landscapes}

\thesisabstract{%
  Machine learning-assisted directed evolution (MLDE) offers a principled framework for navigating
  protein fitness landscapes with limited experimental data.
  A central but underexplored question is whether the encoding--surrogate model pairings
  identified as optimal on dense combinatorial libraries generalise to sparse, full-protein landscapes.
  This thesis systematically benchmarks four encoding strategies ---
  one-hot mutation features, ESM2 mean-pooled embeddings, ESM2 delta embeddings,
  and AAIndex physicochemical parameters --- crossed with three surrogate models ---
  random forests, Gaussian processes, and deep neural network ensembles ---
  on both the four-site GB1 combinatorial library (Wu et al.\ 2019) and the full-protein
  avGFP dataset (Sarkisyan et al.\ 2016).
  We show that ESM2 mean-pooled embeddings yield no significant advantage over one-hot mutation
  features when paired with a random forest surrogate on avGFP ($R^2 = 0.469$ vs.\ $R^2 = 0.466$,
  $n = 10{,}000$), consistent with recent findings on combinatorial landscapes.
  Practical utility is quantified using Precision@K, an operationally meaningful metric measuring
  what fraction of the top-K predicted variants are genuinely high-fitness, which reveals substantially
  lower top-variant recovery than Spearman rank correlation alone implies.
  The results provide actionable guidance for protein engineers choosing ML strategies under
  realistic laboratory budget constraints.
}

\begin{document}
\makecoverpage
\makecopyrightpage
\clearpage

\begin{abstractpage}[english]
Machine learning-assisted directed evolution (MLDE) allows researchers to navigate high-dimensional
protein fitness landscapes more efficiently than classical random mutagenesis by training a
surrogate model on a small labelled set and using it to prioritise the next round of experiments.
A critical but unresolved design question is how to pair the sequence representation
(encoding) with the surrogate model --- and whether the optimal pairing depends on
the statistical structure of the fitness landscape being optimised.

This thesis benchmarks four encoding strategies against three surrogate models on two datasets
with contrasting landscape structure: the four-site GB1 combinatorial library and the full-protein
avGFP dataset from Sarkisyan et al.\ (2016).
We find that ESM2 mean-pooled embeddings provide no significant predictive advantage over sparse
one-hot mutation features when paired with tree-based surrogates, and we quantify practical
top-variant recovery using Precision@K alongside the standard Spearman rank correlation metric.
The results offer actionable guidance for protein engineers choosing ML strategies under
realistic laboratory constraints, and raise the question of whether encoding--surrogate recommendations
from dense combinatorial benchmarks transfer to sparser, full-protein landscapes.
\end{abstractpage}

\newpage
%% Finnish abstract
\thesistitle{Koodausstrategioiden ja surrogaattimallien parituksen vaikutukset koneoppimisavusteisessa proteiinioptimoinnissa}
\thesissubtitle{Systemaattinen vertailu harvoilla ja tihe\"{a}ill\"{a} fitness-maisemilla}
\supervisor{Prof.\ Ohjaajan nimi}
\advisor{Dr Ohjaajan nimi}
\degreeprogram{Elektroniikka ja s\"{a}hk\"{o}tekniikka}
\major{Koneoppiminen ja data-analyysi}
\keywords{proteiinisuunnittelu\spc koneoppiminen\spc bayesilainen optimointi\spc proteiinikielen mallit\spc fitness-maisema}
\begin{abstractpage}[finnish]
Koneoppimisavusteinen ohjattu evoluutio mahdollistaa proteiinien fitness-maisemien
navigoimisen rajoitetulla kokeellisella datalla.
T\"{a}ss\"{a} ty\"{o}ss\"{a} vertaillaan neli\"{a}\"{a} koodausstrategiaa kolmeen surrogaattimalliin
kahdella rakenteeltaan erilaisella datasetill\"{a}: nelipaikkainen GB1-kombinatorinen kirjasto
sek\"{a} Sarkisyanin (2016) avGFP-datasetti.
Tulokset osoittavat, ett\"{a} ESM2-kielimallien mean-pool-upotukset eiv\"{a}t tuota merkitt\"{a}v\"{a}\"{a}
etua yksinkertaisempiin koodausmenetelmiin verrattuna satunnaismets\"{a}surrogaateilla.
Precision@K-metriikka paljastaa merkitt\"{a}v\"{a}sti alhaisemman k\"{a}yt\"{a}nn\"{o}n
huippuvarianttien l\"{o}yt\"{a}miskyvyn kuin Spearman-korrelaatio yksin\"{a}\"{a}n antaisi olettaa.
\end{abstractpage}

\newpage
\dothesispagenumbering{}

\mysection{Preface}
This thesis was carried out at Aalto University during 2025--2026.
I would like to thank my supervisor and advisor for their guidance throughout this project,
and the broader protein engineering and machine learning communities for making the
datasets and tools used here publicly available.

\vspace{2cm}
Otaniemi, \today\\
{\hfill Aleksander Heino \hspace{1cm}}
\newpage

\thesistableofcontents

\mysection{Symbols and abbreviations}
\subsection*{Abbreviations}
\begin{tabular}{ll}
  AAIndex & amino acid index physicochemical parameter database \\
  BO      & Bayesian optimisation \\
  DKL     & deep kernel learning \\
  DMS     & deep mutational scanning \\
  DNN     & deep neural network \\
  EI      & expected improvement \\
  ESM     & evolutionary scale modelling (protein language model family) \\
  FLIP    & fitness landscape inference for proteins (benchmark) \\
  GP      & Gaussian process \\
  LM      & language model \\
  MLDE    & machine learning-assisted directed evolution \\
  P@K     & Precision at K \\
  PLM     & protein language model \\
  RF      & random forest \\
  RBF     & radial basis function (kernel) \\
  UCB     & upper confidence bound \\
  WT      & wild type \\
\end{tabular}

\cleardoublepage

%% ---------------------------------------------------------------
\section{Introduction}
\label{sec:intro}
\thispagestyle{empty}

Proteins are the molecular machines of life.
Whether a protein fluoresces, catalyses a chemical reaction, or binds a target depends
exquisitely on its amino acid sequence, yet only a tiny fraction of all possible sequences
folds into a functional structure.
Directed evolution~\cite{Romero2009, Packer2015} mimics natural selection in the laboratory:
sequences are mutated, expressed, assayed for a target property, and the best-performing
variants are propagated into the next generation of experiments.
While powerful, the approach is inherently combinatorial --- the number of possible sequences grows
as $20^L$ for a protein of length $L$, making exhaustive experimental screening impossible
for all but the shortest peptides.

Machine learning-assisted directed evolution (MLDE)~\cite{Wu2019} was developed as a response
to this combinatorial explosion.
Instead of assaying every possible variant, a small labelled set is used to train a
surrogate fitness model, which predicts fitness across the unseen sequence space and guides
the selection of the next experimental batch.
This closed-loop paradigm --- sometimes framed as Bayesian optimisation (BO) over sequences ---
has been validated on several benchmark datasets and shown to identify high-fitness variants
with substantially fewer experimental measurements than random
screening~\cite{Wu2019, Wittmann2021, Yang2025}.

A central design choice in any MLDE system is how to represent protein sequences to the surrogate
model.
Representations range from simple one-hot encodings of mutated positions to dense embeddings
from protein language models (PLMs) pre-trained on hundreds of millions of evolutionary
sequences~\cite{Lin2023ESM2}.
A recent large-scale benchmark --- Active Learning-assisted Directed Evolution
(ALDE)~\cite{Yang2025} --- established that this choice is not neutral: the optimal encoding
depends on which surrogate model is used.
High-dimensional PLM embeddings benefit from deep learning surrogates; simpler surrogates
such as random forests or Gaussian processes perform best with lower-dimensional one-hot or
physicochemical encodings.

However, ALDE and the majority of MLDE benchmarks were conducted on dense combinatorial libraries
--- datasets in which every combination of mutations at a small number of fixed sites (typically four)
has been measured, yielding on the order of 150,000 variants.
Whether these encoding--surrogate pairing recommendations transfer to sparse, full-protein
landscapes --- where a smaller number of variants carry mutations distributed across the entire
sequence --- is an open question.
The statistical properties of these two landscape types differ qualitatively in ways that
may influence which representation is most informative: in a sparse full-protein landscape,
mean-pooling PLM embeddings over hundreds of positions dilutes the signal from the few mutated
sites, potentially neutralising any advantage of the richer representation.

This thesis investigates that question systematically.
We benchmark four encoding strategies and three surrogate models on two datasets with contrasting
landscape structure: the four-site GB1 combinatorial library~\cite{Wu2019} and the full-protein
avGFP dataset from Sarkisyan et al.~\cite{Sarkisyan2016}, which contains approximately 52,000
variants spanning 1 to 9 mutations across a 238-residue sequence.
In addition to standard metrics (Spearman rank correlation $\rho$, $R^2$), we evaluate
Precision@K --- the fraction of the K highest-predicted variants that genuinely belong to the
top-K by measured fitness --- as an operationally meaningful metric aligned with
the experimental selection bottleneck.

\paragraph{Contributions.}
The main contributions of this thesis are:
\begin{itemize}
  \item a systematic two-dataset benchmark of encoding--surrogate combinations under matched conditions, spanning both dense combinatorial and sparse full-protein landscape regimes;
  \item empirical evidence that ESM2 mean-pooled embeddings provide no significant advantage over one-hot mutation features when paired with a random forest surrogate, on both landscape types;
  \item a characterisation of Precision@K as a metric that reveals substantially lower practical top-variant recovery than Spearman $\rho$ implies, and diverges from $\rho$ across encoding--surrogate combinations; and
  \item an evaluation of delta embeddings (mutant minus wild-type embedding at mutation positions) as a compact, biologically motivated alternative to mean pooling.
\end{itemize}

\paragraph{Thesis structure.}
Section~\ref{sec:litrev} reviews the relevant literature.
Section~\ref{sec:methods} describes the datasets, encodings, surrogates, and evaluation protocol.
Section~\ref{sec:results} presents and analyses the experimental results.
Section~\ref{sec:conclusions} summarises findings, discusses limitations, and outlines
future directions.

\clearpage

%% ---------------------------------------------------------------
\section{Literature Review}
\label{sec:litrev}

\subsection{From Random Mutagenesis to Machine-Assisted Directed Evolution}
\label{sec:litrev:history}

The engineering of proteins for non-natural functions has a history stretching back to the
early 1990s, when Frances Arnold demonstrated that iterative rounds of random mutagenesis and
functional selection could transform the specificity and stability of industrial enzymes
without any knowledge of the underlying molecular mechanism~\cite{Romero2009}.
This paradigm --- \emph{directed evolution} --- exploited natural selection as an optimisation
algorithm: diversity is introduced into a gene pool by error-prone PCR, DNA shuffling, or
saturation mutagenesis~\cite{Stemmer1994}, and the best-performing variants are identified
by a functional assay and carried forward as parents for the next round.
The approach proved remarkably general, and its development was recognised with the Nobel Prize
in Chemistry in 2018.

Classical directed evolution is limited by experimental throughput.
A single round of screening can characterise hundreds to a few thousand variants, but the
sequence space grows combinatorially: for a protein of length $L$ with $k$ variable positions,
there are up to $20^k$ possible sequences.
For the experimentally tractable case of four variable sites, this already yields 160,000
sequences --- at the edge of what a well-equipped laboratory can screen in a single campaign.
For full-protein engineering with mutations distributed across the entire sequence, exhaustive
coverage is impossible.

The introduction of deep mutational scanning (DMS)~\cite{Fowler2014} partially addressed
this bottleneck by coupling high-throughput DNA synthesis and next-generation sequencing to
functional selection.
In a DMS experiment, a library encoding many thousands of protein variants is simultaneously
assayed for function; sequencing before and after selection quantifies the fitness effect of
each variant in a single experiment.
The Sarkisyan et al.\ (2016) study~\cite{Sarkisyan2016} exemplified this approach for the
green fluorescent protein (GFP), characterising approximately 52,000 single-protein mutants
for fluorescence, providing one of the first quantitative maps of a full-protein fitness landscape.
The GB1 dataset from Wu et al.\ (2019)~\cite{Wu2019} extended this to combinatorial libraries,
measuring all $20^4$ combinations at four coevolutionary positions of the protein G B1 domain.

The integration of machine learning into directed evolution --- MLDE~\cite{Wu2019} --- emerged
from the observation that DMS data are rich enough to train a surrogate fitness predictor,
which can then be queried over unevaluated sequences to prioritise the next experimental batch.
Wu et al.\ demonstrated that a random forest regressor trained on a fraction of GB1 variants
could identify top-fitness sequences far more efficiently than random sampling.
Subsequent work formalised the MLDE loop as a Bayesian optimisation problem,
in which the surrogate provides both a fitness estimate and an uncertainty measure,
and an acquisition function trades off exploitation of high-predicted-fitness regions
against exploration of uncertain ones~\cite{Romero2009, Yang2025}.
More recently, the loop has been extended with active learning strategies that jointly optimise
the training set composition and the acquisition policy~\cite{Yang2025},
and with generative models that propose novel sequences outside the training library
rather than selecting from a pre-measured pool~\cite{Stanton2022}.

\subsection{The Machine Learning Toolkit: Encodings, Surrogates, and Acquisition Functions}
\label{sec:litrev:methods}

Three components define an MLDE system: the representation of each sequence (encoding),
the surrogate model that maps representations to predicted fitness,
and the acquisition function that selects which sequences to assay next.

\paragraph{Sequence encodings.}
The simplest representation is a \emph{one-hot} encoding of the amino acid identity at each
variable position, producing a sparse binary vector.
For combinatorial libraries with $k$ variable sites this gives a $20k$-dimensional vector;
for full-protein variants, a sparse vector of mutation identities indexed to the full sequence.
One-hot encodings contain no learned biological information and treat all amino acid substitutions
as equally distant.

Physicochemical encodings such as AAIndex~\cite{Kawashima2000} or the Georgiev reduced
alphabet~\cite{Georgiev2009} map each amino acid to a small number of empirically measured
properties (volume, hydrophobicity, charge, secondary structure propensity), offering
a low-dimensional representation with biological inductive bias.

Protein language models (PLMs) represent the most expressive encoding family currently available.
Trained by masked-token self-supervision on hundreds of millions of evolutionary protein sequences,
models such as ESM1b, ESM2~\cite{Lin2023ESM2}, and ESM3~\cite{Hayes2025ESM3} learn
contextualised, position-specific residue representations that implicitly encode structural and
functional constraints accumulated over evolutionary time.
The ESM2 model family ranges from 8 million to 15 billion parameters; the 650 million parameter
variant (ESM2-650M) produces 1280-dimensional per-position embeddings.
The simplest approach to deriving a fixed-length sequence representation is to average
(\emph{mean-pool}) the per-position embeddings over the full sequence,
yielding a 1280-dimensional vector regardless of sequence length.
An alternative is to extract only the embeddings at variable positions
(\emph{site-specific} extraction), concatenating them to give a higher-dimensional but
mutation-focused vector; this is the approach used in the ALDE benchmark~\cite{Yang2025}.

A third possibility, not yet systematically benchmarked in the literature, is the
\emph{delta embedding}: computing the difference between the mutant and wild-type embeddings
at each mutated position.
This approach is motivated by the observation that mean pooling over a full-length protein
dilutes the mutation signal across hundreds of invariant positions, potentially explaining
why PLM embeddings offer little advantage over one-hot features when paired with simple
surrogate models~\cite{Yang2025}.

\paragraph{Surrogate models.}
The three dominant surrogate families in the MLDE literature are random forests (RF)~\cite{Wu2019},
Gaussian processes (GP)~\cite{Romero2009}, and deep neural network (DNN)
ensembles~\cite{Lakshminarayanan2017}.
Random forests offer simplicity, robustness to small training sets, and out-of-bag uncertainty
estimates via variance across trees, but their uncertainty estimates are known to be poorly
calibrated in sparse input regions.
Gaussian processes provide a principled probabilistic framework with closed-form posterior
uncertainty, and their uncertainty decays correctly in well-explored regions; however,
exact GP inference scales as $\mathcal{O}(n^3)$ in the number of training points, limiting
their use to datasets of a few thousand samples without approximations.
DNN ensembles --- typically five or more independently initialised multi-layer perceptrons ---
offer the expressiveness needed to exploit high-dimensional PLM embeddings and produce
uncertainty estimates via ensemble disagreement~\cite{Lakshminarayanan2017};
they generally require more training data than RF or GP.

Deep kernel learning (DKL)~\cite{Yang2025} combines a DNN feature extractor with a GP covariance
function, aiming to inherit the representational power of deep networks and the calibrated
uncertainty of GPs.
The ALDE benchmark found DKL with ESM2 site-specific embeddings to be one of the top-performing
combinations on GB1 and TrpB, though at substantially higher computational cost than RF or GP.

\paragraph{Encoding-surrogate pairing effects.}
A central finding from Yang et al.~(2025)~\cite{Yang2025} is that the benefit of a richer
encoding depends on which surrogate is used.
High-dimensional PLM embeddings improve performance when paired with DNN surrogates or DKL,
but offer little or no advantage over one-hot or physicochemical encodings when paired with RF or GP.
This \emph{encoding-surrogate pairing} effect was also noted earlier by Wittmann
et al.~\cite{Wittmann2021}, who found that one-hot encoding outperformed complex representations
for non-deep models on GB1.
The practical implication is that investing in expensive PLM featurisation is only justified
when a correspondingly expressive surrogate is used.

\paragraph{Acquisition functions.}
Given a surrogate that outputs a predicted mean $\hat{\mu}$ and uncertainty $\hat{\sigma}$,
an acquisition function decides which sequences to label next.
Greedy selection (rank by $\hat{\mu}$, select top-K) ignores uncertainty but is competitive
in regimes with large training sets.
Upper confidence bound (UCB; select by $\hat{\mu} + \beta\hat{\sigma}$)~\cite{Srinivas2010}
provides a tunable exploration-exploitation trade-off.
Thompson sampling draws a fitness function from the posterior and greedily optimises
it~\cite{Yang2025}, providing implicit diversity in the selected batch.
Expected improvement (EI) maximises the expected gain over the current best observed value.

a non-trivial non-Bayesian baseline is AdaLead~\cite{Sinai2020}, a greedy hill-climbing
algorithm with recombination that does not require a calibrated uncertainty model.
AdaLead has been shown to match or outperform full Bayesian approaches on several benchmarks,
raising the question of when principled uncertainty quantification is worth its additional
complexity.

\paragraph{Zero-shot strategies.}
An orthogonal class of approaches requires no labelled fitness data at all: PLM
log-likelihood ratios between mutant and wild-type sequences are used directly as fitness
proxies~\cite{Notin2022Tranception, Hie2024}.
These zero-shot scores are competitive with supervised surrogates in very-low-data regimes
and provide a strong baseline that any supervised method should exceed before the additional
cost of labelled data is considered justified.
The ProteinGym benchmark~\cite{Notin2023ProteinGym} provides a comprehensive evaluation of
zero-shot predictors across 217 proteins and over 1.5 million variants.

\paragraph{Two paradigms: generative search versus pool-based selection.}
A conceptual divide runs through the protein ML optimisation literature that is rarely stated
explicitly but is critical for interpreting results and assessing prior work.

The \emph{generative search} paradigm --- exemplified by
CbAS~\cite{Brookes2019CbAS}, DbAS~\cite{Brookes2018DbAS},
DynaPPO~\cite{Angermueller2020DynaPPO}, PEX~\cite{Ren2022PEX},
and GGS~\cite{Kirjner2023GGS} --- proposes \emph{novel sequences not present in any
measured dataset}.
A neural network or PLM is trained on available experimental data and then used as an
oracle to score candidate sequences generated by a search algorithm (VAE, RL policy,
Gibbs sampler, greedy hill-climbing).
The research question is which search algorithm most efficiently finds sequences that
the oracle scores as high-fitness.
Surrogate quality is treated as a fixed implementation detail, not a variable.
Importantly, because the oracle is queried on sequences it may never have seen during training,
these methods rely on the oracle's ability to extrapolate --- an assumption that is rarely
validated against experimental ground truth.

The \emph{pool-based MLDE} paradigm --- exemplified by Wu et al.~\cite{Wu2019},
Wittmann et al.~\cite{Wittmann2021}, ALDE~\cite{Yang2025}, and
BOES~\cite{Soldat2025BOES} --- selects the best variants from a \emph{fixed pool of
already-measured sequences}.
A small subset of the pool (the training budget, typically $n = 96$--$5000$) is used to
train a surrogate; the surrogate ranks the remaining pool members; the top-ranked
variants are selected.
Ground truth is always a real experimental measurement, not a neural network
approximation.
The research question is which encoding and surrogate combination best recovers the
highest-fitness variants from the pool given the available training budget.

A third category --- \emph{low-N generative design}~\cite{Biswas2021LowN} ---
combines limited training labels ($n = 24$--$96$) with generative search over novel
sequences and experimental wet-lab validation of top candidates.
This is the closest in spirit to real-world directed evolution campaigns, but it does not
directly address the encoding-surrogate pairing question for pool-based prediction.

This thesis belongs to the pool-based MLDE paradigm.
Generative search papers frequently use the avGFP dataset but do so by training an
oracle on all available measurements and proposing novel sequences --- a different
task that does not evaluate surrogate quality under a limited labelling budget.
Therefore, the extensive use of avGFP in the generative literature does not constitute
prior work addressing the research question investigated here.

\subsection{Benchmark Datasets and Fitness Landscape Structure}
\label{sec:litrev:datasets}

Progress in MLDE has been anchored by a small number of landmark datasets.
Understanding their structural properties --- not just their size --- is essential for
interpreting benchmark results and for assessing whether findings generalise across settings.

\paragraph{GB1 (Wu et al., 2019).}
The GB1 dataset~\cite{Wu2019} is the de facto standard benchmark for combinatorial MLDE.
It covers all $20^4 = 160{,}000$ amino acid combinations at four coevolving positions
(V39, D40, G41, V54) in the B1 domain of streptococcal protein G,
of which 149,361 variants were successfully measured for binding to IgG-Fc by
deep mutational scanning.
The landscape is highly epistatic: the fitness effect of any single mutation depends strongly
on the identities of the other three positions.
Approximately 92\% of measured variants have near-zero fitness (the ``hole'' problem;
Wittmann et al.\ 2021~\cite{Wittmann2021}), making the training distribution heavily skewed
and posing a challenge for supervised models trained on random subsets.

\paragraph{TrpB (Johnston et al., 2024).}
The TrpB dataset~\cite{Johnston2024} covers four sites in the indole-3-glycerol phosphate
synthase subunit TrpB, providing a second dense combinatorial benchmark with a different
epistatic structure from GB1.
Johnston et al.\ characterised approximately 160,000 four-site variants for thermostability
and activity, enabling the ALDE benchmark~\cite{Yang2025} to test whether GB1 findings
generalise across combinatorial landscapes.

\paragraph{avGFP (Sarkisyan et al., 2016).}
The avGFP dataset~\cite{Sarkisyan2016} represents a fundamentally different landscape structure
from the combinatorial benchmarks above.
Rather than exhaustively measuring all combinations at a few fixed sites, Sarkisyan et al.\
characterised approximately 52,000 random mutants of the Aequorea victoria green fluorescent
protein (avGFP), each carrying between one and nine mutations distributed across the full
238-residue sequence.
The resulting landscape is \emph{sparse} --- it covers only a tiny fraction of the sequence
space reachable by a small number of mutations --- and heterogeneous in mutation count
and location.

The avGFP dataset has been used in both paradigms described above.
In the search paradigm, a CNN or similar model is trained on all 52,000 measurements and
used as an oracle; the research question is which search algorithm finds sequences the oracle
predicts as highly fluorescent.
In the surrogate paradigm relevant to this thesis, only a small subset of measurements is
available as training data, and the question is which encoding-surrogate combination best
ranks the remaining variants.
The latter has received substantially less attention despite being the operationally relevant
setting for wet-lab protein engineering.

Several structural features distinguish avGFP from combinatorial benchmarks.
First, the training distribution is not a random subset of a combinatorial grid but a biased
sample of an enormous space; models must generalise to sequences with different mutation counts
and positions simultaneously.
Second, the fraction of near-zero-fitness variants (hole density) is lower than in GB1,
partially reducing the class imbalance problem.
Third, mean-pooling PLM embeddings over the full 238-residue sequence dilutes the mutation
signal across more than 230 invariant positions, a structural artefact that may explain why
PLM featurisation offers limited gains over one-hot features when the surrogate is a random
forest (see Section~\ref{sec:results:esm2-onehot}).
Fourth, ALDE's site-specific embedding --- which extracts embeddings at a fixed small set of
variable positions --- does not directly apply to avGFP, where variable positions differ
from variant to variant.

\paragraph{ProteinGym and FLIP.}
The ProteinGym benchmark~\cite{Notin2023ProteinGym} aggregates DMS data for 217 proteins and
evaluates zero-shot fitness predictors.
The FLIP benchmark~\cite{Dallago2021FLIP} provides standardised splits for supervised fitness
prediction.
Both benchmarks are valuable for broad coverage, but neither addresses the encoding-surrogate
pairing question directly, and neither disentangles the dense-combinatorial versus sparse-full-protein
landscape regimes.

\paragraph{Evaluation metrics.}
Spearman's rank correlation coefficient $\rho$ is the dominant evaluation metric in both
zero-shot benchmarking~\cite{Notin2023ProteinGym} and MLDE studies~\cite{Wu2019, Yang2025}.
It measures the monotonic association between predicted and measured fitness ranks across the
full test set.
The $R^2$ coefficient of determination is also widely reported.

However, in the context of active sequence selection, what matters most operationally is not
global rank correlation but the quality of the top of the predicted distribution:
a model may achieve high $\rho$ overall yet fail to concentrate the genuinely best variants
at the top of its ranked list.
Precision@K --- the fraction of the K highest-predicted variants that actually belong to the
top-K by measured fitness --- captures this directly and aligns naturally with the experimental
bottleneck of a fixed synthesis-and-assay budget.
Despite this operational relevance, Precision@K is underreported in the published MLDE
literature; most papers cite Spearman $\rho$ or $R^2$ as their primary metrics.

\subsection{Current State of the Art and Open Questions}
\label{sec:litrev:sota}

The ALDE benchmark~\cite{Yang2025} is the most systematic published comparison of
encoding-surrogate-acquisition function combinations to date.
Testing all combinations of four encodings (one-hot, AAIndex/Georgiev, ESM2 site-specific),
four surrogate types (boosting ensemble, GP, DNN ensemble, DKL), and two acquisition functions
(greedy, Thompson sampling) on GB1 and TrpB, it established that:
(i)\ the DNN ensemble with Thompson sampling is the best overall strategy across conditions;
(ii)\ ESM2 site-specific embeddings improve performance when paired with DNN or DKL surrogates;
and (iii)\ non-deep surrogates benefit little or not at all from richer encodings.
These findings are consistent with the earlier observation by Wittmann et al.~\cite{Wittmann2021}
that one-hot encoding outperforms more complex representations for non-deep models on GB1.

EVOLVEpro~\cite{Liu2024EVOLVEpro} demonstrated a practical active learning system using
ESM2-15B mean-pooled embeddings with a random forest surrogate, achieving measurable fitness
improvements across five diverse protein targets in an iterative wet-lab setting.
EVOLVEpro also reported a notable failure mode: PLM zero-shot fitness scores (log-likelihood
ratios) showed near-zero or negative correlation with experimentally measured activity for
several proteins, highlighting that evolutionary fitness and specific biological function are
distinct quantities.

ESM3~\cite{Hayes2025ESM3} introduced a multimodal generative PLM jointly modelling sequence,
structure, and function tokens.
At 98 billion parameters, the full model is computationally prohibitive for most research
groups; however, smaller ESM C variants (300 million parameters) are more accessible and may
provide embeddings competitive with ESM2 at lower computational cost.
Whether ESM3 embeddings provide material benefits over ESM2 in supervised fitness prediction
relevant to MLDE is not yet established.

Three important questions remain unresolved at the time of writing.
First, essentially all systematic encoding-surrogate benchmarks have been conducted on dense
combinatorial libraries (GB1, TrpB); whether optimal pairings transfer to sparse full-protein
landscapes such as avGFP is unknown.
Second, the delta embedding approach has been proposed as a biologically motivated alternative
to mean pooling but has not been benchmarked.
Third, Precision@K as a primary evaluation metric for active learning in protein engineering
remains underused; it is not clear how strongly encoding-surrogate combinations differ on
this metric versus Spearman $\rho$.

\subsection{Research Question and Contributions}
\label{sec:litrev:rq}

The preceding review motivates the central research question of this thesis:

\begin{quote}
  \emph{Do the encoding-surrogate pairing recommendations derived from dense combinatorial
  fitness landscapes (GB1) transfer to sparse, full-protein landscapes (avGFP),
  and how does the choice of encoding-surrogate pair affect Precision@K as an
  operationally meaningful evaluation metric?}
\end{quote}

More concretely, we test whether the relative ranking of four encoding strategies
(one-hot mutation features, ESM2 mean-pooled, ESM2 delta, AAIndex physicochemical)
across three surrogate models (random forest, Gaussian process, DNN ensemble) is consistent
between GB1 and avGFP, and whether Spearman $\rho$ and Precision@K give concordant or
divergent rankings.

The work makes four concrete contributions.
First, it provides the first systematic encoding-surrogate benchmark that spans both the
dense combinatorial and sparse full-protein landscape regimes.
Second, it evaluates Precision@K alongside Spearman $\rho$ and $R^2$, quantifying the
gap between global correlation and practical top-variant recovery.
Third, it tests the delta embedding hypothesis: that subtracting wild-type embeddings at
mutation positions recovers signal lost by mean pooling, potentially restoring a PLM advantage
in the sparse regime.
Fourth, it characterises learning curves for all combinations, providing guidance on the
minimum labelling budget at which each encoding-surrogate pair is competitive.

\clearpage

%% ---------------------------------------------------------------
\section{Methods}
\label{sec:methods}
\thispagestyle{empty}

\subsection{Datasets}
\label{sec:methods:data}

\paragraph{GB1.}
The GB1 dataset is obtained from the SaProtHub repository (SaProtHub/Dataset-GB1-fitness)
and the FLIP benchmark~\cite{Dallago2021FLIP}, using the standard two-versus-rest split
in which any variant containing at least one of the four wild-type amino acids is held out
for testing.
% TODO: confirm exact split and filtering; describe hole handling (near-zero threshold)

\paragraph{avGFP.}
The avGFP dataset is obtained from the MaveDB repository using the Sarkisyan et al.\
accession (MaveDB:00000055-a-1) and is preprocessed to retain only
variants with valid fluorescence measurements and between one and nine amino acid substitutions
relative to the wild type.
% TODO: confirm exact filtering; report final n_train and n_test

\subsection{Sequence Encodings}
\label{sec:methods:encodings}

\paragraph{One-hot mutation features.}
Each variant is represented by a sparse binary vector of length proportional to the number
of variable positions, encoding the identity of the amino acid change at each mutated site.
For GB1 this yields a $4 \times 20 = 80$-dimensional vector; for avGFP the vector length
varies with mutation count and is zero-padded to a fixed maximum.
% TODO: describe exact padding strategy for variable-length avGFP variants

\paragraph{ESM2 mean-pooled embeddings.}
Each variant sequence is passed through ESM2-650M (Meta AI; Lin et al.\ 2023~\cite{Lin2023ESM2})
and the per-position embeddings are averaged over all residues to produce a 1280-dimensional
sequence representation.
Embeddings are precomputed and stored; ESM2 is not fine-tuned.

\paragraph{ESM2 delta embeddings.}
The mutant embedding is computed by subtracting the wild-type embedding at each mutated
position: $\delta_i = h^\text{mut}_i - h^\text{WT}_i$ for each mutated position $i$.
The per-position deltas are averaged to produce a 1280-dimensional summary vector.
This approach is designed to concentrate the representation on mutation-induced changes
and reduce the dilution effect of averaging over invariant positions.
% TODO: confirm implementation; handle multi-site variants

\paragraph{AAIndex physicochemical features.}
Each amino acid substitution is encoded using a selected subset of AAIndex~\cite{Kawashima2000}
indices covering hydrophobicity, volume, charge, and secondary structure propensity.
% TODO: specify which indices; describe aggregation for multi-site variants

\subsection{Surrogate Models}
\label{sec:methods:surrogates}

\paragraph{Random forest.}
A random forest regressor is trained using scikit-learn with $n_\text{estimators} = 100$
and default hyperparameters.
% TODO: report full hyperparameter set; describe out-of-bag uncertainty estimate

\paragraph{Gaussian process.}
An exact GP with a radial basis function (RBF) kernel is trained using GPyTorch~\cite{Gardner2018}.
Kernel hyperparameters are optimised by marginal likelihood maximisation.
Inputs are standardised to zero mean and unit variance before fitting.
% TODO: describe kernel; training set size cap for avGFP (GP scales as O(n^3))

\paragraph{DNN ensemble.}
An ensemble of five independent multi-layer perceptrons is trained with identical architecture
but different random initialisations~\cite{Lakshminarayanan2017}.
Each network has two hidden layers of 256 units with ReLU activations.
Uncertainty is estimated as the standard deviation of ensemble predictions.
% TODO: confirm architecture; training details (lr, batch size, epochs)

\subsection{Evaluation Protocol}
\label{sec:methods:eval}

All models are evaluated in an \emph{offline} Bayesian optimisation setting:
a fixed training set of size $n$ is sampled uniformly at random from the full dataset,
and the surrogate is used to rank the remaining (unseen) variants.
Each configuration is repeated across five independent random seeds.

Three metrics are reported: Spearman rank correlation $\rho$ between predicted and
measured fitness on the held-out set; $R^2$ on the same set; and Precision@K for
$K \in \{10, 100, 1000\}$.
Precision@K is defined as
\begin{equation}
  \text{P@}K = \frac{|\{k : \text{rank}_\text{pred}(k) \leq K\} \cap \{k : \text{rank}_\text{true}(k) \leq K\}|}{K},
\end{equation}
where $\text{rank}_\text{pred}(k)$ is the rank of variant $k$ in the predicted ordering and
$\text{rank}_\text{true}(k)$ is its rank by measured fitness.

Learning curves are generated by varying training set size
$n \in \{50, 100, 200, 500, 1000, 2000, 5000, 10000\}$ with five seeds per configuration.

\clearpage

%% ---------------------------------------------------------------
\section{Results}
\label{sec:results}
\thispagestyle{empty}

\subsection{ESM2 Mean Pooling versus One-Hot Features on avGFP}
\label{sec:results:esm2-onehot}

\paragraph{Preliminary finding (avGFP, RF, n=10{,}000).}
A random forest surrogate trained on ESM2 mean-pooled embeddings (1280-dimensional, ESM2-650M)
achieves $R^2 = 0.469 \pm 0.003$ on the avGFP held-out set,
virtually identical to the same surrogate trained on one-hot mutation features
($R^2 = 0.466 \pm 0.001$, $n = 5$ independent runs).
The difference is not statistically significant.
Despite comparable $R^2$, Precision@100 is low for both encoding strategies (3--9\%),
indicating that the models correctly rank the broad fitness distribution but struggle to
concentrate the top-100 highest-fitness variants at the top of their predicted lists.

% TODO: add full results table across all encoding x surrogate x dataset combinations
% TODO: add Spearman rho values and learning curve figure

\subsection{Learning Curves}
\label{sec:results:learning-curves}

% TODO: report crossover training set size at which ESM2 and one-hot strategies diverge
% (preliminary data suggests crossover around n~200-500 for Spearman rho on avGFP)

\subsection{GB1 Benchmark}
\label{sec:results:gb1}

% TODO: report GB1 results once experiments are run
% Expected: results should be consistent with Yang et al. (2025) ALDE findings

\subsection{Delta Embedding Evaluation}
\label{sec:results:delta}

% TODO: compare delta embedding vs mean pool vs one-hot; test hypothesis that delta
% embedding partially recovers PLM advantage in sparse regime

\clearpage

%% ---------------------------------------------------------------
\section{Discussion and Conclusions}
\label{sec:conclusions}
\thispagestyle{empty}

\subsection{Summary of Findings}

% TODO: summarise main results once experiments are complete

\subsection{Implications for Protein Engineering Practice}

% TODO: practical guidance for choosing encoding-surrogate combinations under different budgets
% Key message: for RF/GP surrogates, one-hot or physicochemical encodings are competitive with ESM2
% and much cheaper to compute. ESM2 embeddings are only worth the compute cost when paired with
% a DNN or DKL surrogate and sufficient training data.

\subsection{Limitations}

% TODO: discuss offline-only evaluation (no iterative wet-lab loop), single-protein-per-regime
% limitation, absence of multi-round active learning simulation

\subsection{Future Work}

Promising directions include:
replacing the random forest with an AdaLead greedy baseline to test whether BO-style
acquisition adds value over simple hill-climbing;
extending the benchmark to the TrpB dataset to confirm GB1 findings generalise
across dense combinatorial landscapes;
evaluating ESM3 embeddings as replacements for ESM2;
and simulating multi-round active learning with observation noise to approximate the
online setting studied in the wet-lab MLDE literature.

\clearpage
\printbibliography[heading=bibintoc]

%% Optional appendix
\clearpage
\thesisappendix
\section{Experimental Configuration}
\label{app:config}

% TODO: full hyperparameter tables, software versions, compute environment details

\end{document}
